% appendix-d-implementation-index.tex -- part of thesis_v3.tex; do not add \documentclass here.
% Generated from the register pass on Chapters 2 and 3. Each row records the repository
% artifact that a statement in the running text was taken from.

\chapter{Implementation Index}
\label{ch:appendix_implementation_index}

Chapters~\ref{ch:theory_background} and~\ref{ch:design_implementation} describe the method and
the system in the usual narrative form, without interrupting the text with source-file references.
This appendix restores that traceability. Each row names a statement made in the running text and
the artifact in the project repository it was taken from, grouped by the section in which the
statement appears. Line numbers are omitted deliberately: they move as the code changes, whereas
file and symbol names are stable. Paths are relative to the repository root.

\fillin{state the repository URL and the commit hash the thesis was written against, so a reader can resolve these paths to an exact revision}

\section{Chapter 2 --- Theoretical Background}
\label{sec:appD_ch2}

\subsection*{Impedance Behavior and Decoupling Strategies}
\begin{longtable}{p{0.56\textwidth} p{0.36\textwidth}}
\hline
Statement in the text & Source artifact \\
\hline
\endfirsthead
\hline
Statement in the text & Source artifact \\
\hline
\endhead
Anti-resonance enters the surrogate as a shape prior: every impedance peak must be preceded by a resonance dip, scored as mean(peak\_score * exp(-10 * prior\_dip)) on the raw log-impedance channel. & \texttt{experiments/\allowbreak exp059\_capacity\_freq/\allowbreak codes/\allowbreak physics\_loss.py-233} \\
From bin index 190 upwards the spectrum is monotonically increasing by construction, so the anti-resonance term acts mainly at low frequency. & \texttt{experiments/\allowbreak exp059\_capacity\_freq/\allowbreak codes/\allowbreak physics\_loss.py-233 (docstring)} \\
The target mask is a fixed array of 231 values, minimum 0.8 and maximum 49.99999999999999, stored with shape (231,1) and squeezed to (231,); Stage 2 aborts if the squeezed shape is not (231,). & \texttt{configs/\allowbreak target\_impedance.npy; pipelines/\allowbreak latent/\allowbreak optimize.py-548} \\
\hline
\end{longtable}

\subsection*{Power Integrity Simulation and Multi-modal Design Data}
\begin{longtable}{p{0.56\textwidth} p{0.36\textwidth}}
\hline
Statement in the text & Source artifact \\
\hline
\endfirsthead
\hline
Statement in the text & Source artifact \\
\hline
\endhead
Six-step data flow: sample/merge combination CSVs, build PEB project in ECADSTAR script order, simulate via headless engineer.exe --batch on Windows, append to the raw multi-frequency dataset, robust log1p per-MHz normalization with optional K<=30 filter, then train/evaluate/active-learn. & \texttt{docs/\allowbreak data-pipeline.md, section 1} \\
Row i of a combination CSV maps to the ECADSTAR object PI-(i+1) in output order. & \texttt{docs/\allowbreak data-pipeline.md, section 1} \\
The conventional manual design loop the framework replaces: place, simulate spectrum, on mask violation run a spatial PI distribution, move capacitors toward hotspots, repeat. & \texttt{docs/\allowbreak framework-overview.md, section 1} \\
Board carries 52 decap slots on a 7x8 grid with four invalid cells at (0,3), (0,4), (6,3), (6,4); ordered labels C1-C52, index i of b maps to LABELS\_ORDERED[i]. & \texttt{libs/\allowbreak data\_creation/\allowbreak occupancy.py-31} \\
Slot count is derived from the label list: NUM\_SLOTS = len(LABELS\_ORDERED). & \texttt{experiments/\allowbreak exp059\_capacity\_freq/\allowbreak codes/\allowbreak graph\_occ.py} \\
Occupancy tensors are loaded from the layout store, flattened, clamped to [0,1], and the budget recovered as K = sum\_i b\_i. & \texttt{src\_vae/\allowbreak others/\allowbreak dataloader.py-201} \\
The spectrum has exactly 231 points; the raw CSV reader rejects any file of different length (EXPECTED\_IMP\_LENGTH = 231). & \texttt{libs/\allowbreak data\_creation/\allowbreak impedance.py,33-36} \\
Spectrum batches of shape (B,231) are unsqueezed to (B,1,231) inside the model. & \texttt{experiments/\allowbreak exp059\_capacity\_freq/\allowbreak codes/\allowbreak train\_core.py-736} \\
The impedance decoder returns shape (B,1,231). & \texttt{experiments/\allowbreak exp059\_capacity\_freq/\allowbreak codes/\allowbreak vae\_multi\_input\_simple.py} \\
The frequency grid is a stored array of shape (231,) running 1000000.0 Hz, 1100000.0 Hz, ... 600000000.0 Hz; first log10 step 0.0414, last 0.0073; no specification document defines it. & \texttt{configs/\allowbreak Frequency\_data\_hz.npy} \\
The PI-distribution heatmap is interpolated onto a 64x64 grid (HEATMAP\_GRID\_SIZE = 64); raw creation returns the two-channel stack [Zi * mask\_board, mask\_board]. & \texttt{libs/\allowbreak data\_creation/\allowbreak heatmap.py,74} \\
Normalized heatmaps that reach training are single-channel float32 of shape (1,64,64). & \texttt{datasets/\allowbreak data\_multifreq\_train\_norm\_unbounded/\allowbreak heatmap/\allowbreak sample\_1000000.npy} \\
Heatmap encoder first layer is nn.Conv2d(1, 16, ...), consistent with single-channel input. & \texttt{experiments/\allowbreak exp059\_capacity\_freq/\allowbreak codes/\allowbreak vae\_multi\_input\_simple.py} \\
Heatmap decoder ends in nn.Conv2d(16, 1, kernel\_size=3, padding=1), i.e. one channel at 64x64. & \texttt{experiments/\allowbreak exp059\_capacity\_freq/\allowbreak codes/\allowbreak vae\_multi\_input\_simple.py} \\
The training set declares 24 PI frequency anchors in MHz (10.0 ... 500.0). & \texttt{datasets/\allowbreak data\_multifreq\_train\_norm\_unbounded/\allowbreak dataset\_meta.json, key pi\_frequencies\_mhz} \\
The same 24-anchor list appears in the normalization statistics under Heatmap.anchors\_mhz and PI\_freq.anchor\_mhz. & \texttt{datasets/\allowbreak data\_multifreq\_train\_norm\_unbounded/\allowbreak normalization\_stats.json} \\
Anchors are resolved in order from dataset\_meta.json, the environment data directory, a legacy manifest, configs/multifreq\_anchors.yaml, then a built-in default tuple (10,63,80,130,150,200,250,270,330,400,500); the YAML file lists 20 entries. Both fall-backs disagree with the live 24-anchor dataset and are marked legacy. & \texttt{src\_vae/\allowbreak others/\allowbreak multifreq\_anchors.py-21,49-81; configs/\allowbreak multifreq\_anchors.yaml} \\
PI\_freq normalization constants for the current dataset: log10\_min = 6.0, log10\_range = 2.778151250383644. & \texttt{datasets/\allowbreak data\_multifreq\_train\_norm\_unbounded/\allowbreak normalization\_stats.json, key PI\_freq} \\
When a batch carries no frequency the training code substitutes the normalized value for 200.0 MHz, the default\_mhz of pi\_freq\_norm\_for\_model. & \texttt{src\_vae/\allowbreak others/\allowbreak multifreq\_anchors.py (pi\_freq\_norm\_for\_model)} \\
Normalized training dataset: 582997 manifest rows over 24979 unique layouts, layout store, max\_k\_filter 30; per-anchor counts 24979 except 80 MHz with 8480. & \texttt{datasets/\allowbreak data\_multifreq\_train\_norm\_unbounded/\allowbreak dataset\_meta.json} \\
The K <= 30 filter drops layouts with K > 30 before statistics and output are computed. & \texttt{docs/\allowbreak dataset.md, section 3} \\
Heatmap normalization robust\_log1p\_per\_mhz\_unbounded with per-MHz median/IQR, background -3.1792, clip\_min -1.6791870594024658, clip\_max 6.275787830352783; the same three values appear in the exp059 config. & \texttt{datasets/\allowbreak data\_multifreq\_train\_norm\_unbounded/\allowbreak normalization\_stats.json; experiments/\allowbreak exp059\_capacity\_freq/\allowbreak config.yaml-142} \\
The unbounded loader path applies no clip. & \texttt{src\_vae/\allowbreak others/\allowbreak dataloader.py-217} \\
Impedance normalization is a log z-score with mu\_log = -1.039238452911377 and sigma\_log = 1.7831873893737793, fitted over 24979 layouts. & \texttt{datasets/\allowbreak data\_multifreq\_train\_norm\_unbounded/\allowbreak normalization\_stats.json} \\
Layout-level holdout: all frequencies of one design\_id go to train or validation but never both; train\_split 0.9, seed 42. & \texttt{experiments/\allowbreak exp059\_capacity\_freq/\allowbreak config.yaml,19; experiments/\allowbreak exp059\_capacity\_freq/\allowbreak codes/\allowbreak dataloader\_multifreq.py-80} \\
\hline
\end{longtable}

\subsection*{Formulation of the Decoupling Capacitor Placement Problem}
\begin{longtable}{p{0.56\textwidth} p{0.36\textwidth}}
\hline
Statement in the text & Source artifact \\
\hline
\endfirsthead
\hline
Statement in the text & Source artifact \\
\hline
\endhead
Engineering problem statement: keep Z(f) below the target mask over 1-600 MHz at 231 bins with b in \{0,1\}\textasciicircum{}52; exhaustive search over 2\textasciicircum{}52 placements is intractable. & \texttt{docs/\allowbreak framework-overview.md, section 1} \\
Feasibility is tested against the mask shifted inward by BOUNDARY\_MARGIN = 0.1. & \texttt{pipelines/\allowbreak latent/\allowbreak optimize.py,357-363} \\
The default comparison space is LOSS\_SPACE = 'log', i.e. z-scored log-impedance, so the margin 0.1 is in normalized units and not 0.1 ohm. & \texttt{pipelines/\allowbreak latent/\allowbreak optimize.py} \\
Per-bin violation exponent EXCEED\_POWER = 2.0; base objective weights w\_gap = 1.0 and w\_exc = 25.0; default OBJECTIVE\_MODE = 'gap\_max', alternative 'feasible\_map' keeps only the exceed term. & \texttt{pipelines/\allowbreak latent/\allowbreak optimize.py,100-102,313-329} \\
Selection metric defaults to SELECT\_METRIC = 'max\_ohm'; the best candidate is only updated for feasible seeds, and a no\_solution.json is written per K when no seed is feasible. & \texttt{pipelines/\allowbreak latent/\allowbreak optimize.py,366-369,466-471,644-646} \\
Stage 2 loads the exp038\_true\_multi VAE and impedance surrogate (EXPERIMENT = 'exp038\_true\_multi'); no Stage-2 entry point is wired to exp059\_capacity\_freq. & \texttt{pipelines/\allowbreak latent/\allowbreak optimize.py,45-47} \\
NORMALIZATION\_STATS\_PATH points at datasets/data\_multifreq\_norm/normalization\_stats.json, a directory absent from the current checkout; the live dataset is datasets/data\_multifreq\_train\_norm\_unbounded. & \texttt{pipelines/\allowbreak latent/\allowbreak optimize.py} \\
\hline
\end{longtable}

\subsection*{Variational Autoencoders and Multi-modal Learning}
\begin{longtable}{p{0.56\textwidth} p{0.36\textwidth}}
\hline
Statement in the text & Source artifact \\
\hline
\endfirsthead
\hline
Statement in the text & Source artifact \\
\hline
\endhead
exp059 KL settings: free\_bits 0.12, beta annealing off with beta\_final 0.03, per-expert KL weight 0.02, sigma-regularization weight 1.5 toward target 0.45. & \texttt{experiments/\allowbreak exp059\_capacity\_freq/\allowbreak config.yaml-118} \\
PoE fusion: precision-weighted product of Gaussian experts, a unit-Gaussian prior expert always appended, eps = 1e-8 added to each variance and to the combined precision. & \texttt{experiments/\allowbreak exp059\_capacity\_freq/\allowbreak codes/\allowbreak vae\_multi\_input\_simple.py-449} \\
Expert log-variances clamped to [-6.0, 1.0] at each head; fused log-variance clamped to [-4.0, 2.0] before reparameterization. & \texttt{experiments/\allowbreak exp059\_capacity\_freq/\allowbreak codes/\allowbreak vae\_poe\_freq.py,160-162} \\
Latent partition arithmetic: d\_shared = d\_latent - d\_hm\_priv - d\_imp\_peak and d\_layout = d\_shared + d\_imp\_peak. & \texttt{experiments/\allowbreak exp059\_capacity\_freq/\allowbreak codes/\allowbreak vae\_multi\_input\_simple.py-133,171-174} \\
exp059 capacity: latent\_dim 128, heatmap\_private\_dim 40, cond\_dim 32. & \texttt{experiments/\allowbreak exp059\_capacity\_freq/\allowbreak config.yaml-4} \\
imp\_peak\_dim is unset in the exp059 config and takes the builder default 8, alongside graph\_hidden\_dim 128, graph\_num\_layers 3, graph\_dropout 0.1. & \texttt{experiments/\allowbreak exp059\_capacity\_freq/\allowbreak codes/\allowbreak exp059\_common.py-54} \\
d\_shared = 80 is independently stated in the experiment notes (d\_layout = 88 follows from the partition arithmetic). & \texttt{experiments/\allowbreak exp059\_capacity\_freq/\allowbreak notes.md} \\
Expert-to-latent wiring: heatmap expert writes the full latent, occupancy expert the shared dims, impedance expert shared+peak; unowned dims padded with the prior. Heatmap encoder input width 64 + 2*cond\_dim; occupancy/impedance heads 64 + cond\_dim. & \texttt{experiments/\allowbreak exp059\_capacity\_freq/\allowbreak codes/\allowbreak vae\_multi\_input\_simple.py-258,394-413} \\
A dedicated frequency expert maps the frequency embedding to the heatmap-private dims only; use\_freq\_poe\_expert defaults to True and is not overridden in the exp059 config. & \texttt{experiments/\allowbreak exp059\_capacity\_freq/\allowbreak codes/\allowbreak vae\_poe\_freq.py-84; experiments/\allowbreak exp059\_capacity\_freq/\allowbreak codes/\allowbreak exp059\_common.py} \\
Occupancy slot graph: 52 nodes, four-neighbour edges on the 7x8 grid, self-loops, symmetric normalization D\textasciicircum{}-1/2 A D\textasciicircum{}-1/2; node features occupancy (1) + slot embedding (hidden/2) + normalized row/col (2), positions normalized by 6.0 and 7.0; mean+max pooling projected to width 64. & \texttt{experiments/\allowbreak exp059\_capacity\_freq/\allowbreak codes/\allowbreak graph\_occ.py-47} \\
Spectrum chain graph over 231 bins with optional self-loops and the same symmetric normalization; node features bin value (1) + bin embedding (hidden/2) + normalized position (1) from torch.linspace(0.0, 1.0, 231), pooled to width 64. & \texttt{experiments/\allowbreak exp059\_capacity\_freq/\allowbreak codes/\allowbreak graph\_imp.py,13-24} \\
Graph convolution layer computes lin\_self(x) + lin\_neigh(Ax) followed by LayerNorm, LeakyReLU(0.1) and optional dropout, in plain PyTorch with no torch\_geometric dependency. & \texttt{experiments/\allowbreak exp059\_capacity\_freq/\allowbreak codes/\allowbreak graph\_occ.py-75 (imported by graph\_imp.py)} \\
The two 64-wide graph features are concatenated and projected by Linear(128, 64) with LayerNorm and LeakyReLU into the joint feature feeding the occupancy and impedance expert heads. & \texttt{experiments/\allowbreak exp059\_capacity\_freq/\allowbreak codes/\allowbreak vae\_multi\_input\_simple.py-247,422-428} \\
Graph decoders broadcast latent+conditioning to every node and read out 52 occupancy logits / 231 spectrum values. & \texttt{experiments/\allowbreak exp059\_capacity\_freq/\allowbreak codes/\allowbreak graph\_occ.py-178; experiments/\allowbreak exp059\_capacity\_freq/\allowbreak codes/\allowbreak graph\_imp.py-123} \\
Budget embedding table has 53 entries covering K in \{0,...,52\}; the heatmap conditioning vector is the concatenation of the K and frequency embeddings, width 2*cond\_dim. & \texttt{experiments/\allowbreak exp059\_capacity\_freq/\allowbreak codes/\allowbreak vae\_multi\_input\_simple.py-193,457-468} \\
Frequency embedding is a Fourier-feature MLP over [f, sin(2*pi*f*b), cos(2*pi*f*b)] for b = 1..B, input width 1 + 2B, mapped Linear-SiLU-Linear-SiLU to cond\_dim; exp059 sets freq\_fourier\_features 16 (exp057/exp058 set 8). The exp059 model imports this module from exp043. & \texttt{experiments/\allowbreak exp043/\allowbreak codes/\allowbreak freq\_conditioning.py-42} \\
FiLM applies x -> x*(1+gamma) + beta with gamma/beta from a Linear layer whose weights and bias are zero-initialized except that the bias entries of the gamma half are set to 1.0. & \texttt{experiments/\allowbreak exp043/\allowbreak codes/\allowbreak freq\_conditioning.py-59} \\
Multi-scale FiLM at four decoder scales: 8x8 on 128 channels, 16x16 on 64, 32x32 on 32, 64x64 on 16, each conditioned on the 2*cond\_dim vector; a further FiLM2d(32, 2*cond\_dim) sits in the heatmap encoder. & \texttt{experiments/\allowbreak exp059\_capacity\_freq/\allowbreak codes/\allowbreak vae\_multi\_input\_simple.py-375,505-544} \\
use\_multiscale\_film is absent from the exp059 config and defaults to True in the model-kwargs builder; the exp057 builder has no such key. & \texttt{experiments/\allowbreak exp059\_capacity\_freq/\allowbreak codes/\allowbreak exp059\_common.py} \\
The exp059 heatmap decoder has no U-Net skip connections: latent plus FiLM only, to match Stage 2 and generation. exp057's builder still exposes use\_heatmap\_unet\_skips defaulting to True. & \texttt{experiments/\allowbreak exp059\_capacity\_freq/\allowbreak codes/\allowbreak vae\_multi\_input\_simple.py} \\
Modality dropout: inside encode(), each of the three data experts is kept with probability 1 - modality\_dropout via torch.bernoulli(torch.full((3,), 1.0 - modality\_dropout)); if all three drop, one is forced back on; the frequency expert is appended afterwards and never dropped; modality\_dropout\_protect\_heatmap forces the heatmap expert on. & \texttt{experiments/\allowbreak exp059\_capacity\_freq/\allowbreak codes/\allowbreak vae\_poe\_freq.py-158} \\
exp059 modality-dropout values: 0.06 with the same start value, heatmap protection from epoch 25 / fraction 0.325, focus-phase 0.03. & \texttt{experiments/\allowbreak exp059\_capacity\_freq/\allowbreak config.yaml-33,46,118-119} \\
encode\_occupancy\_latent\_full builds a zero impedance batch, runs the joint layout encoder and forms the posterior from the occupancy expert plus the PoE prior alone; its docstring records that this matches the AL inference path and the occupancy-only fine-tune. & \texttt{experiments/\allowbreak exp059\_capacity\_freq/\allowbreak codes/\allowbreak vae\_multi\_input\_simple.py-905} \\
With probability layout\_train\_prob the decode latent comes from the layout branch, and within it with probability occ\_only\_encode\_prob from the occupancy-only encode; the full encode is always run for the KL term. & \texttt{experiments/\allowbreak exp059\_capacity\_freq/\allowbreak codes/\allowbreak train\_core.py-715} \\
Config-vs-notes disagreement: live config sets layout\_train\_prob 0.65 and occ\_only\_encode\_prob 0.4 and distillation weights 2.5/2.5; the AL fine-tune override sets 0.9 and 0.7; the notes record 0/0 and 0/0 plus different eval\_use\_occ\_only\_layout, al\_overlay\_data\_dir, learning\_rate, num\_epochs and resume\_checkpoint. & \texttt{experiments/\allowbreak exp059\_capacity\_freq/\allowbreak config.yaml,48,159,162; experiments/\allowbreak exp059\_capacity\_freq/\allowbreak config\_al\_finetune.yaml; experiments/\allowbreak exp059\_capacity\_freq/\allowbreak notes.md; CLAUDE.md} \\
Layout-to-private head: two-hidden-layer MLP over [occ\_feat(64), imp\_feat(64), cond(2*cond\_dim)] predicting the heatmap-private dims, optional frequency FiLM, squashed as 4*tanh(x/4); acts as distillation student for the heatmap encode teacher. use\_layout\_private\_head and use\_layout\_private\_freq\_film default True and are unset in the exp059 config. Weights: latent 2.5, private multiplier 3.0, logvar 0.3, output 2.5. & \texttt{experiments/\allowbreak exp059\_capacity\_freq/\allowbreak codes/\allowbreak vae\_multi\_input\_simple.py-291; experiments/\allowbreak exp059\_capacity\_freq/\allowbreak codes/\allowbreak vae\_poe\_freq.py-264} \\
\hline
\end{longtable}

\subsection*{Physics-Informed and Surrogate Modeling}
\begin{longtable}{p{0.56\textwidth} p{0.36\textwidth}}
\hline
Statement in the text & Source artifact \\
\hline
\endfirsthead
\hline
Statement in the text & Source artifact \\
\hline
\endhead
PhysicsCritic CNN maps occ\_spatial (B,1,7,8) plus PI\_freq to an effectiveness map (B,1,64,64) in [0,1], supervised by inverting and min-max normalizing the heatmap foreground, supplying a radius-influence penalty; detached during the VAE step. & \texttt{experiments/\allowbreak exp059\_capacity\_freq/\allowbreak codes/\allowbreak physics\_loss.py-52,76,164-217} \\
Physics weights: radius-influence 1.0, critic supervision 2.0, anti-resonance 0.5. & \texttt{experiments/\allowbreak exp059\_capacity\_freq/\allowbreak config.yaml-126} \\
Occupancy is binarized to a hard top-K vector before the heatmap decode via occupancy\_for\_heatmap\_decode(occupancy, K, force\_binary=True, ste=self.occupancy\_binary\_ste), vectorized through a descending-sort threshold so per-row variable K needs no Python loop. & \texttt{experiments/\allowbreak exp059\_capacity\_freq/\allowbreak codes/\allowbreak vae\_multi\_input\_simple.py-646; experiments/\allowbreak exp059\_capacity\_freq/\allowbreak codes/\allowbreak occupancy\_binary.py-65} \\
exp059 sets occupancy\_binary\_decode = true and occupancy\_binary\_ste\_train = false. & \texttt{experiments/\allowbreak exp059\_capacity\_freq/\allowbreak config.yaml-165} \\
Decoder input widths: heatmap decoder takes full latent + 2*cond\_dim; occupancy decoder z\_shared + K embedding; impedance decoder z\_layout + K embedding + 64-wide occupancy graph context (passed occupancy, else sigmoid(occ\_recon)). & \texttt{experiments/\allowbreak exp059\_capacity\_freq/\allowbreak codes/\allowbreak vae\_multi\_input\_simple.py-296,625-654} \\
The frequency-weighted 2D-FFT heatmap loss (L1 between log1p magnitudes of rfft2, optional Gaussian MHz bump) is implemented but disabled: heatmap\_spectral\_weight is absent from the config and from the Config dataclass, so the lookup falls back to 0.0 and the function returns None. The exp059 notes claim weight 1.5, mid-boost 1.5, centre 200 MHz, width 80 MHz. & \texttt{experiments/\allowbreak exp059\_capacity\_freq/\allowbreak codes/\allowbreak train\_vae\_simple.py-222,199-201; experiments/\allowbreak exp059\_capacity\_freq/\allowbreak notes.md} \\
eval\_off\_anchor\_mhz = [155.0, 250.0, 265.0] with weights 2.5, 3.0, 2.5; 250.0 MHz is itself one of the 24 training anchors. & \texttt{experiments/\allowbreak exp059\_capacity\_freq/\allowbreak config.yaml-60; datasets/\allowbreak data\_multifreq\_train\_norm\_unbounded/\allowbreak dataset\_meta.json} \\
decode\_heatmap\_blended log-blends the decodes at the two bracketing anchors with (lo, hi, w) from bracket\_anchors\_mhz(mhz); the training and evaluation path decodes directly at the requested PI\_freq. & \texttt{experiments/\allowbreak exp059\_capacity\_freq/\allowbreak codes/\allowbreak vae\_multi\_input\_simple.py-694} \\
AL candidate generation settings: K from 3 to 10, 1600 candidates per K, 80 worst per K, 640 ECAD simulations per cycle, heatmap grid 64, evaluation MHz grid [120,150,155,170,180,200,230,250,265,270,280,300]. & \texttt{active\_learning\_pi/\allowbreak config/\allowbreak exp059\_gp\_error.json} \\
\hline
\end{longtable}

\subsection*{Optimization in Latent Representation Spaces}
\begin{longtable}{p{0.56\textwidth} p{0.36\textwidth}}
\hline
Statement in the text & Source artifact \\
\hline
\endfirsthead
\hline
Statement in the text & Source artifact \\
\hline
\endhead
Straight-through top-K read-out implemented as hard + occ\_prob - occ\_prob.detach(); the docstring records that without it the top-K argsort severs the gradient and z never moves in response to the spectrum target. & \texttt{pipelines/\allowbreak latent/\allowbreak optimize.py-268} \\
Stage-2 search settings: K\_LIST = 1..25, NUM\_CANDIDATE\_SEEDS 32, NUM\_STEPS 1200, LR 5e-2, GRAD\_CLIP 5.0. & \texttt{pipelines/\allowbreak latent/\allowbreak optimize.py-57,71-73} \\
Total Stage-2 objective weights (z-L2 0.05, shape 2.0, posterior 10.0, track 2.0, anti-resonance 0.5, peak w\_peak, diversity 0.35, occupancy-confidence 0.5) and POSTERIOR\_SIGMA\_LIMIT = 2.5. & \texttt{pipelines/\allowbreak latent/\allowbreak optimize.py-97,372-416} \\
The impedance path defaults to a separate occupancy-to-spectrum surrogate: USE\_SURROGATE = \_env\_flag('LATENT\_OPT\_USE\_SURROGATE', '1'). & \texttt{pipelines/\allowbreak latent/\allowbreak optimize.py,271-282} \\
That surrogate maps occupancy (52,) to a log-z spectrum (1,231) with hidden widths (256, 512, 512, 512, 256); its docstring states PI\_freq is not used. & \texttt{experiments/\allowbreak exp038\_true\_multi/\allowbreak codes/\allowbreak surrogate\_impedance.py-8,36-53} \\
\hline
\end{longtable}

\subsection*{Research Gap and Motivation for the Proposed Framework}
\begin{longtable}{p{0.56\textwidth} p{0.36\textwidth}}
\hline
Statement in the text & Source artifact \\
\hline
\endfirsthead
\hline
Statement in the text & Source artifact \\
\hline
\endhead
Retargeting to a new Z\_target(f) by re-running the latent search with the decoder frozen is a property of the implemented Stage-2 script, which is wired to exp038\_true\_multi. & \texttt{pipelines/\allowbreak latent/\allowbreak optimize.py,45-47} \\
\hline
\end{longtable}

\section{Chapter 3 --- Design and Implementation}
\label{sec:appD_ch3}

\subsection*{Raw Power-Integrity Simulation Outputs and File Structure}
\begin{longtable}{p{0.56\textwidth} p{0.36\textwidth}}
\hline
Statement in the text & Source artifact \\
\hline
\endfirsthead
\hline
Statement in the text & Source artifact \\
\hline
\endhead
The raw ECADSTAR export is converted into a layout store of layouts/, heatmap/ and PI\_freq/ trees plus a manifest.csv index. & \texttt{pipelines/\allowbreak data/\allowbreak processing\_multifreq.py} \\
The processing stage follows the repository convention of a module-level CONFIGURATION block instead of command-line arguments. & \texttt{pipelines/\allowbreak data/\allowbreak processing\_multifreq.py} \\
Per design the store keeps occ.npy of shape (52,) and imp.npy of shape (231,); the occupancy encoder is sized by NUM\_SLOTS = len(LABELS\_ORDERED) and the impedance encoder by NUM\_BINS = 231. & \texttt{pipelines/\allowbreak data/\allowbreak processing\_multifreq.py} \\
Only the active-learning batch invocation of the simulator is documented (engineer.exe <design.erf> --batch <file.peb> --batch-auto-exit); bulk data generation is not described anywhere. & \texttt{docs/\allowbreak data-pipeline.md sec.5} \\
The 64x64 heatmap resolution is inferred from the AL config heatmap\_grid\_size = 64 and the decoder's final FiLM stage, not read from dataset metadata. & \texttt{active\_learning\_pi/\allowbreak config/\allowbreak exp059\_gp\_error.json} \\
\hline
\end{longtable}

\subsection*{Multi-frequency Dataset Construction and Normalization}
\begin{longtable}{p{0.56\textwidth} p{0.36\textwidth}}
\hline
Statement in the text & Source artifact \\
\hline
\endfirsthead
\hline
Statement in the text & Source artifact \\
\hline
\endhead
Dataset stage 'normalized', source script pipelines/normalize/multifreq.py, timestamp 2026-07-06T13:39:03.768907+00:00; 582997 manifest rows, 24979 unique layouts, 24979 layout dirs, 582997 heatmap and 582997 PI\_freq files; 24 anchor frequencies; samples\_per\_mhz 24979 for all anchors except 80 MHz at 8480. & \texttt{datasets/\allowbreak data\_multifreq\_train\_norm\_unbounded/\allowbreak dataset\_meta.json} \\
Heatmap normalization mode robust\_log1p\_per\_mhz\_unbounded with unbounded = true; per-MHz medians and IQRs (10.0 MHz: 0.05887173116207123 / 0.03866304084658623; 500.0 MHz: 1.3578120470046997 / 0.6529430150985718); clip\_min -1.6791870594024658, clip\_max 6.275787830352783, z\_max 11.23223876953125, background\_value -3.1792. Impedance log\_mean -1.039238452911377, log\_std 1.7831873893737793, layout\_count 24979, z\_min -2.7170498371124268, z\_max 2.474900722503662. Frequency log10\_min 6.0, log10\_range 2.778151250383644, min\_hz 1000000.0, max\_hz 600000000.0. & \texttt{datasets/\allowbreak data\_multifreq\_train\_norm\_unbounded/\allowbreak normalization\_stats.json} \\
Foreground/background threshold BG\_OHM = 0.05; robust clip percentiles 0.5 and 99.5 apply only in bounded mode; env switches NORM\_ROBUST\_PER\_MHZ (default '1') and NORM\_UNBOUNDED\_Z (default '0'). & \texttt{pipelines/\allowbreak normalize/\allowbreak multifreq.py} \\
The unbounded build filters layouts at MAX\_K = 30, with K recovered per design as (occ > 0.5).sum(), dropping them from both the statistics fit and the written output. & \texttt{pipelines/\allowbreak normalize/\allowbreak build\_train\_norm\_unbounded.py} \\
The experiment config mirrors the heatmap normalization constants (background\_value -3.1792, heatmap\_z\_clip\_min -1.6791870594024658, heatmap\_z\_clip\_max 6.275787830352783) and sets heatmap\_clip\_recon = false. & \texttt{experiments/\allowbreak exp059\_capacity\_freq/\allowbreak config.yaml} \\
Documentation conflict on the anchor set: a 'Full 16-MHz set' is claimed here. & \texttt{docs/\allowbreak dataset.md sec.4} \\
Documentation conflict on the anchor set: anchors named as (10, 70, 120, 200, 270, 400) MHz here. & \texttt{docs/\allowbreak normalization-and-losses.md sec.1.4} \\
\hline
\end{longtable}

\subsection*{Data Splitting, Sampling, and Batch Construction}
\begin{longtable}{p{0.56\textwidth} p{0.36\textwidth}}
\hline
Statement in the text & Source artifact \\
\hline
\endfirsthead
\hline
Statement in the text & Source artifact \\
\hline
\endhead
Layout-level holdout: split\_by\_design = true, train\_split = 0.9; \_split\_indices\_by\_design assigns all rows of a design\_id to one side via torch.randperm; default split seed 42. The docstring records that an earlier frequency-level split 'leaked layout identity into val and is unsuitable for GP / AL residual evaluation'. & \texttt{experiments/\allowbreak exp059\_capacity\_freq/\allowbreak codes/\allowbreak dataloader\_base.py} \\
stratify\_by\_k only enables precompute\_k on the dataset; it does not stratify the split. & \texttt{experiments/\allowbreak exp059\_capacity\_freq/\allowbreak codes/\allowbreak dataloader\_base.py} \\
The weighted sampler is constructed only when balance\_k or balance\_freq is true; both are false in exp059, so no K-stratification or K/frequency balancing is active. & \texttt{experiments/\allowbreak exp059\_capacity\_freq/\allowbreak codes/\allowbreak dataloader\_base.py-362} \\
Config sets stratify\_by\_k = true (conflicting with the dataloader behaviour), balance\_k = false, balance\_freq = false. & \texttt{experiments/\allowbreak exp059\_capacity\_freq/\allowbreak config.yaml} \\
Layout holdout dated 2026-07-20; seed 42 and train\_split 0.9 give approximately 22481 train and 2498 validation layouts (stated with a tilde); required a fresh retrain with resume\_checkpoint null, prior artifacts archived under archive\_pre\_layout\_holdout\_*. & \texttt{experiments/\allowbreak exp059\_capacity\_freq/\allowbreak notes.md-62} \\
train\_samples\_per\_epoch = 64000 with RandomSampler with replacement; batch\_size 224, val\_batch\_size 320, num\_workers 4, prefetch\_factor 2, persistent\_workers false, cache\_in\_ram false. & \texttt{experiments/\allowbreak exp059\_capacity\_freq/\allowbreak config.yaml} \\
Cross-frequency pair lookup maps each training index to alternate-MHz rows of the same design; make\_multifreq\_collate\_with\_pairs attaches the paired row to the batch. & \texttt{experiments/\allowbreak exp059\_capacity\_freq/\allowbreak codes/\allowbreak dataloader\_base.py} \\
\hline
\end{longtable}

\subsection*{Architectural Design}
\begin{longtable}{p{0.56\textwidth} p{0.36\textwidth}}
\hline
Statement in the text & Source artifact \\
\hline
\endfirsthead
\hline
Statement in the text & Source artifact \\
\hline
\endhead
Decoder input widths: dec\_in\_occ = shared\_latent\_dim + cond\_dim; dec\_in\_imp = layout\_latent\_dim + cond\_dim + 64 ('z\_layout + K + occ graph ctx'); heatmap decoder takes full z plus 2*cond\_dim. & \texttt{experiments/\allowbreak exp059\_capacity\_freq/\allowbreak codes/\allowbreak vae\_multi\_input\_simple.py-296} \\
layout\_latent\_dim = shared\_latent\_dim + imp\_peak\_dim. & \texttt{experiments/\allowbreak exp059\_capacity\_freq/\allowbreak codes/\allowbreak vae\_multi\_input\_simple.py} \\
The shared latent block is 80 dimensions for the exp059 configuration. & \texttt{experiments/\allowbreak exp059\_capacity\_freq/\allowbreak notes.md} \\
FiLM2d and FreqConditioner are imported from an earlier fork, so the exp059 stack is not self-contained. & \texttt{experiments/\allowbreak exp059\_capacity\_freq/\allowbreak codes/\allowbreak vae\_multi\_input\_simple.py -> experiments/\allowbreak exp043/\allowbreak codes/\allowbreak freq\_conditioning.py} \\
U-Net skip modules (OptionalSkipFuse, \_last\_heatmap\_skips, heatmap\_skips) were deleted; decoder is latent+FiLM only. & \texttt{experiments/\allowbreak exp059\_capacity\_freq/\allowbreak notes.md-51} \\
use\_multiscale\_film is absent from config.yaml and defaults to True in vae\_model\_kwargs, so the headline architectural change of exp059 is governed by a code default. & \texttt{experiments/\allowbreak exp059\_capacity\_freq/\allowbreak codes/\allowbreak exp059\_common.py} \\
latent\_dim 128, heatmap\_private\_dim 40, cond\_dim 32, freq\_fourier\_features 16, use\_heatmap\_film true; imp\_peak\_dim unset (constructor default 8); graph hyperparameters unset (defaults graph\_hidden\_dim 128, graph\_num\_layers 3, graph\_dropout 0.1). & \texttt{experiments/\allowbreak exp059\_capacity\_freq/\allowbreak config.yaml} \\
Modality dropout schedule: modality\_dropout 0.06, modality\_dropout\_start 0.06, modality\_dropout\_anneal\_epochs 18, modality\_dropout\_protect\_heatmap\_epoch 25, heatmap\_focus\_modality\_dropout 0.03; occupancy\_binary\_decode true, occupancy\_binary\_ste\_train false. & \texttt{experiments/\allowbreak exp059\_capacity\_freq/\allowbreak config.yaml} \\
Model class MultiInputVAEPoeFreq: PoE fusion subclass adding a PI\_freq expert on the private dims and removing U-Net skips; expert logvar clamp [-6.0, 1.0], fused PoE logvar clamp [-4.0, 2.0]; encoders (52-slot 7x8 graph, 231-bin chain graph, CNN with HM\_BOTTLENECK 8, HM\_DROPOUT 0.05); layout\_joint\_fusion Linear(128,64)+LayerNorm+LeakyReLU(0.1); nn.Embedding(53, cond\_dim); layout\_private\_hidden 384 with 4.0*tanh(mu/4.0) squash; topk\_occ\_binary vectorised read-out; encode paths encode, encode\_cross\_modal, encode\_layout\_latent\_full, encode\_occupancy\_latent. & \texttt{experiments/\allowbreak exp059\_capacity\_freq/\allowbreak codes/\allowbreak vae\_multi\_input\_simple.py} \\
\hline
\end{longtable}

\subsection*{Objective Functions and Physics-Informed Regularization}
\begin{longtable}{p{0.56\textwidth} p{0.36\textwidth}}
\hline
Statement in the text & Source artifact \\
\hline
\endfirsthead
\hline
Statement in the text & Source artifact \\
\hline
\endhead
The 2D-FFT log-magnitude L1 anti-blur loss (\_heatmap\_spectral\_loss) is implemented, with an optional Gaussian MHz bump. & \texttt{experiments/\allowbreak exp059\_capacity\_freq/\allowbreak codes/\allowbreak train\_vae\_simple.py-221} \\
The heatmap\_spectral\_weight lookup returns the default 0.0 because the key is absent from config.yaml, and the function returns None for any non-positive weight, so the term never enters the total. & \texttt{experiments/\allowbreak exp059\_capacity\_freq/\allowbreak codes/\allowbreak train\_vae\_simple.py} \\
The experiment notes present the frequency-weighted spectral loss as one of three headline exp059 changes, with heatmap\_spectral\_weight 1.5, mid\_boost 1.5, centre 200 MHz, width 80 MHz -- keys that do not exist in config.yaml. & \texttt{experiments/\allowbreak exp059\_capacity\_freq/\allowbreak notes.md-23} \\
PhysicsLoss is instantiated in the training core, and its parameters reach the optimizer and the checkpoints. & \texttt{experiments/\allowbreak exp059\_capacity\_freq/\allowbreak codes/\allowbreak train\_core.py} \\
The trainer monkey-patches the live epoch function (\_tr.\_run\_epoch = \_run\_epoch\_encode). & \texttt{experiments/\allowbreak exp059\_capacity\_freq/\allowbreak codes/\allowbreak train\_vae\_simple.py} \\
The replacement epoch function drops the physics and pw arguments from its signature. & \texttt{experiments/\allowbreak exp059\_capacity\_freq/\allowbreak codes/\allowbreak run\_epoch\_encode.py} \\
The replacement epoch function calls the loss with physics=None. & \texttt{experiments/\allowbreak exp059\_capacity\_freq/\allowbreak codes/\allowbreak run\_epoch\_encode.py} \\
\_vae\_loss\_exp055 likewise forwards physics=None. & \texttt{experiments/\allowbreak exp059\_capacity\_freq/\allowbreak codes/\allowbreak train\_vae\_simple.py} \\
train\_core.vae\_loss guards the radius-influence, critic-supervision and anti-resonance terms with 'if physics is not None', so none of them is added in exp059. & \texttt{experiments/\allowbreak exp059\_capacity\_freq/\allowbreak codes/\allowbreak train\_core.py-613} \\
Physics module implements radius influence, critic supervision (PhysicsCritic, dual-branch, freq\_emb\_dim 16, detached inside the radius term) and the anti-resonance formulation of eq:anti\_resonance. & \texttt{experiments/\allowbreak exp059\_capacity\_freq/\allowbreak codes/\allowbreak physics\_loss.py} \\
Physics weights configured non-zero: physics\_ri\_weight 1.0, physics\_critic\_sup\_weight 2.0, physics\_ar\_weight 0.5, physics\_critic\_warmup\_epochs 11, physics\_slope\_anneal\_epochs 22, physics\_fg\_clip\_min -1.6791870594024658. & \texttt{experiments/\allowbreak exp059\_capacity\_freq/\allowbreak config.yaml} \\
Reconstruction weights heatmap 5.0, occupancy 7.0, impedance 4.0, occ\_k\_consistency 2.0, with heatmap-focus overrides from epoch 25 (heatmap 7.0, impedance 4.0, cross-freq 1.1, dynamic range 3.0); occupancy focal gamma 1.5 with pos-weight (52-K)/K clamped [0.5, 10.0]; impedance composite weights (deriv 1.4, topk\_k 20, topk 4.5, under 2.8, concavity 2.5, freq alpha 2.0, dual topk 0.75, peak index 1.5, peak mag 2.0, num\_peaks 8, peak start 25 / ramp 22 / focus 25); heatmap shape terms (pearson 0.75, grad vector 2.0, grad direction 1.0, min grad magnitude 0.08, grad huber delta 0.5, top-region 2.5 at q=0.95 and 2.5 at q=0.05 with delta 0.5, peak-loc 2.5 at temperature 0.035, valley-loc 0.0); latent terms (free\_bits 0.12, use\_beta\_annealing false, beta\_final 0.03, mu\_hinge\_weight 0.1 threshold 4.0, mu\_bias\_weight 0.05, per\_expert\_kl\_weight 0.02, sigma\_reg\_target 0.45, sigma\_reg\_weight 1.5, upper band 1.3x target); cross-frequency (weight 1.2, start epoch 20, layout\_mix\_prob 0.45, use\_encode\_z true); distillation (latent 2.5 with private x3.0 and logvar 0.3, output 2.5 Huber vs detached teacher); frequency jitter (prob 0.6, log10 sigma 0.1). & \texttt{experiments/\allowbreak exp059\_capacity\_freq/\allowbreak config.yaml} \\
\hline
\end{longtable}

\subsection*{Training, Inference}
\begin{longtable}{p{0.56\textwidth} p{0.36\textwidth}}
\hline
Statement in the text & Source artifact \\
\hline
\endfirsthead
\hline
Statement in the text & Source artifact \\
\hline
\endhead
Training is launched as a module: 'python -m experiments.exp059\_capacity\_freq.codes.train\_vae\_simple' with VAE\_EXPERIMENT\_DIR set and VAE\_CONFIG\_PATH unset; two-GPU runs use 'torchrun --standalone --nproc\_per\_node=2 -m experiments.exp059\_capacity\_freq.codes.train\_vae\_simple'. (Invocation mechanics; removed from the body.) & \texttt{experiments/\allowbreak exp059\_capacity\_freq/\allowbreak codes/\allowbreak train\_vae\_simple.py} \\
load\_yaml\_config strips '\#' comment lines and calls json.loads, so config.yaml is JSON despite its extension; an optional VAE\_CONFIG\_PATH environment override is layered on top of config.yaml. (File-format mechanics; removed from the body, override mechanism retained.) & \texttt{experiments/\allowbreak exp059\_capacity\_freq/\allowbreak codes/\allowbreak exp059\_common.py} \\
The exp059 trainer monkey-patches train\_core: Config, build\_vae\_model, \_run\_epoch, heatmap\_loss, vae\_loss, \_on\_stats\_loaded, the dataloader factory, the YAML loader and the impedance loss. & \texttt{experiments/\allowbreak exp059\_capacity\_freq/\allowbreak codes/\allowbreak train\_vae\_simple.py} \\
AdamW with weight\_decay 1e-4 (fused, then foreach, then plain fallback on CUDA); ReduceLROnPlateau min mode with lr\_factor 0.5, lr\_patience 15, lr\_min 1e-5; impedance\_separate\_lr true, impedance\_decoder\_lr\_mult 1.0, impedance\_peak\_focus\_epoch 25. & \texttt{experiments/\allowbreak exp059\_capacity\_freq/\allowbreak config.yaml} \\
The experiment notes state the exp059 fresh run used learning\_rate 2e-4, num\_epochs 500 and resume\_checkpoint null. & \texttt{experiments/\allowbreak exp059\_capacity\_freq/\allowbreak notes.md} \\
The current config states learning\_rate 1.5e-05, num\_epochs 510, reset\_lr\_on\_resume true and resume\_checkpoint experiments/exp059\_capacity\_freq/checkpoints/last\_model.pt -- no longer matching the notes. & \texttt{experiments/\allowbreak exp059\_capacity\_freq/\allowbreak config.yaml,8,138} \\
Numerical-stability policy: amp 'off', tf32 false, compile false, training\_force\_fp32 true, training\_skip\_nonfinite\_batches true, training\_grad\_clip\_norm 1.0, logvar clamp [-10.0, 10.0], recon z soft floor -8.0, finite loss cap 1000000.0; clipping applied as torch.nn.utils.clip\_grad\_norm\_(params, 1.0). & \texttt{experiments/\allowbreak exp059\_capacity\_freq/\allowbreak config.yaml, experiments/\allowbreak exp059\_capacity\_freq/\allowbreak codes/\allowbreak train\_core.py} \\
Curriculum milestones: cross\_freq\_start\_epoch 20, heatmap\_focus\_start\_epoch 25, modality\_dropout\_protect\_heatmap\_epoch 25, impedance\_peak\_start\_epoch 25, beta\_end\_epoch 25, physics\_critic\_warmup\_epochs 11, penalty\_warmup\_epochs 11. & \texttt{experiments/\allowbreak exp059\_capacity\_freq/\allowbreak config.yaml} \\
Validation on epoch 1, each checkpoint interval and the final epoch (val\_on\_checkpoint\_only true), checkpoint\_interval 5, keep\_last\_n\_checkpoints 0; off-anchor evaluation every 5 epochs at 155.0/250.0/265.0 MHz with weights 2.5/3.0/2.5, max 20 batches, eval\_use\_occ\_only\_layout true. & \texttt{experiments/\allowbreak exp059\_capacity\_freq/\allowbreak config.yaml} \\
The off-anchor metric doubles as a checkpoint-selection criterion. & \texttt{experiments/\allowbreak exp059\_capacity\_freq/\allowbreak checkpoints/\allowbreak best\_off\_anchor\_model.pt} \\
Each run writes loss.csv, heatmap\_peak\_split.csv, impedance\_split.csv, off\_anchor\_eval.csv, epoch\_timing.csv and timing.json. & \texttt{experiments/\allowbreak exp059\_capacity\_freq/\allowbreak metrics/\allowbreak } \\
DDP enabled with ddp\_base\_batch\_size 160 and ddp\_linear\_lr\_scale true; nccl backend with a default 30-minute timeout overridable by NCCL\_TIMEOUT\_MINUTES; model wrapped with find\_unused\_parameters=True. & \texttt{experiments/\allowbreak exp059\_capacity\_freq/\allowbreak codes/\allowbreak distributed\_train.py, experiments/\allowbreak exp059\_capacity\_freq/\allowbreak config.yaml} \\
Encode-path mixing: layout\_train\_prob 0.65 and occ\_only\_encode\_prob 0.4. & \texttt{experiments/\allowbreak exp059\_capacity\_freq/\allowbreak config.yaml-48} \\
The notes describe the exp059 capacity phase as full-encode only, with both mixing probabilities at 0. & \texttt{experiments/\allowbreak exp059\_capacity\_freq/\allowbreak notes.md-14} \\
The handover document likewise describes the capacity phase as full-encode only (probabilities 0). & \texttt{docs/\allowbreak cursor-handoff.md-37} \\
The model parameter count is printed at runtime but stored in no artifact. & \texttt{experiments/\allowbreak exp059\_capacity\_freq/\allowbreak codes/\allowbreak train\_core.py} \\
\hline
\end{longtable}

\subsection*{Auxiliary Guidance Network for Latent-Space Optimization}
\begin{longtable}{p{0.56\textwidth} p{0.36\textwidth}}
\hline
Statement in the text & Source artifact \\
\hline
\endfirsthead
\hline
Statement in the text & Source artifact \\
\hline
\endhead
The nearest artifact to an 'auxiliary guidance network' is a feed-forward occupancy-to-impedance surrogate: occupancy (52,) to impedance log-z (1, 231), PI\_freq unused; hidden\_dims (256, 512, 512, 512, 256), dropout 0.1, num\_epochs 200, batch\_size 128, learning\_rate 3e-4, lr\_patience 15, lr\_factor 0.5, lr\_min 1e-5, weight\_decay 1e-4, topk\_k 20, topk\_weight 7.0, under\_penalty 2.8. Present only in exp037\_lat\_change, exp038\_true\_multi, exp039\_improved\_heatmap, exp040 and exp041; absent from exp057, exp058, exp059, exp060. & \texttt{experiments/\allowbreak exp038\_true\_multi/\allowbreak codes/\allowbreak surrogate\_impedance.py} \\
Stage 2 loads the impedance surrogate checkpoint when USE\_SURROGATE is on (the default), and otherwise falls back to the VAE impedance decoder head. & \texttt{experiments/\allowbreak exp038\_true\_multi/\allowbreak checkpoints/\allowbreak surrogate\_best.pt} \\
\hline
\end{longtable}

\subsection*{Problem Formulation}
\begin{longtable}{p{0.56\textwidth} p{0.36\textwidth}}
\hline
Statement in the text & Source artifact \\
\hline
\endfirsthead
\hline
Statement in the text & Source artifact \\
\hline
\endhead
Stage 2 defaults to EXPERIMENT = 'exp038\_true\_multi' for both the VAE occupancy decoder and the impedance surrogate checkpoint, and imports MultiInputVAE from that fork; feasibility margin BOUNDARY\_MARGIN 0.1, objective mode 'gap\_max', LOSS\_SPACE 'log', EXCEED\_POWER 2.0, gap weight 1.0, exceed weight 25.0; target mask loaded from configs/target\_impedance.npy, required shape (231,) after squeeze, converted by (log Z - mu\_log)/sigma\_log. & \texttt{pipelines/\allowbreak latent/\allowbreak optimize.py} \\
Exact form of the exceed maximum-excursion component and of the gap term (not extracted; subject of a \textbackslash\{\}fillin). & \texttt{pipelines/\allowbreak latent/\allowbreak optimize.py-329} \\
Stage-2 default normalization statistics path is datasets/data\_multifreq\_norm/normalization\_stats.json, a directory absent from the current checkout. & \texttt{pipelines/\allowbreak latent/\allowbreak optimize.py} \\
The exp059 training set is datasets/data\_multifreq\_train\_norm\_unbounded, which is the path Stage 2 should use. & \texttt{experiments/\allowbreak exp059\_capacity\_freq/\allowbreak config.yaml} \\
The limitation that Stage 2 is not wired to exp059 is also recorded in the project documentation. & \texttt{docs/\allowbreak latent-optimization.md sec.7} \\
The stored target mask has shape (231, 1), dtype float64, minimum 0.8 and maximum 49.99999999999999, with no unit metadata and no located generating script. & \texttt{configs/\allowbreak target\_impedance.npy} \\
\hline
\end{longtable}

\subsection*{Optimization Procedure and Feasibility Assessment}
\begin{longtable}{p{0.56\textwidth} p{0.36\textwidth}}
\hline
Statement in the text & Source artifact \\
\hline
\endfirsthead
\hline
Statement in the text & Source artifact \\
\hline
\endhead
Search settings: K\_LIST 1..25, NUM\_CANDIDATE\_SEEDS 32, NUM\_STEPS 1200, LR 5e-2, GRAD\_CLIP 5.0, INIT\_SHARED\_TEMP 0.8, OPTIMIZE\_BATCH\_PER\_K True; seeds random per run (saved to run\_config.json) or fixed 0..N-1 via LATENT\_OPT\_SEED\_MODE, default 'random'. & \texttt{pipelines/\allowbreak latent/\allowbreak optimize.py} \\
Straight-through top-K (\_ste\_topk) with the in-code justification 'Without this the top-K argsort severs the gradient and z never moves on the spectrum target'; applied only on the surrogate branch, surrogate(\_ste\_topk(occ, K)), while the fallback model.decode\_impedance(z, K) branch keeps the read-out out of the gradient path. USE\_SURROGATE defaults to True. & \texttt{pipelines/\allowbreak latent/\allowbreak optimize.py} \\
Search regularizers: Z\_L2\_WEIGHT 0.05 with Z\_PRIOR\_MODE 'agg\_posterior', POSTERIOR\_BOUNDARY\_WEIGHT 10.0 beyond POSTERIOR\_SIGMA\_LIMIT 2.5, SHAPE\_REG\_WEIGHT 2.0 with SHAPE\_MIN\_STD\_DLOG 0.12, SHAPE\_TRACK\_WEIGHT 2.0, PEAK\_LOSS\_WEIGHT 1.0, PEAK\_INDEX\_WEIGHT 3.0, PEAK\_MAG\_WEIGHT 2.0, DUAL\_TOPK\_WEIGHT 2.5, IMPEDANCE\_NUM\_PEAKS 8, IMPEDANCE\_TOPK\_K 20 (all overridable via LATENT\_OPT\_* environment variables), DIVERSITY\_WEIGHT 0.35, OCC\_CONFIDENCE\_WEIGHT 0.5, PHYSICS\_AR\_WEIGHT 0.5 overridable by the checkpoint config's physics\_ar\_weight. & \texttt{pipelines/\allowbreak latent/\allowbreak optimize.py} \\
Selection metric SELECT\_METRIC = 'max\_ohm' (alternative 'mean'); infeasible budgets emit an explicit no\_solution.json rather than a best-effort placement. & \texttt{pipelines/\allowbreak latent/\allowbreak optimize.py} \\
\hline
\end{longtable}

\subsection*{Experimental Setup}
\begin{longtable}{p{0.56\textwidth} p{0.36\textwidth}}
\hline
Statement in the text & Source artifact \\
\hline
\endfirsthead
\hline
Statement in the text & Source artifact \\
\hline
\endhead
ECADSTAR host settings: erf\_path C:\textbackslash\{\}Users\textbackslash\{\}muthusamy\textbackslash\{\}Desktop\textbackslash\{\}design\textbackslash\{\}H-shape.emc\textbackslash\{\}H-shape.erf, powerbus Power\_GND, peb\_components IC1\_Port1, heatmap\_only\_peb true, wait\_timeout\_sec 18000. (Absolute Windows path removed from the body.) & \texttt{active\_learning\_pi/\allowbreak config/\allowbreak exp059\_gp\_error.json} \\
ECADSTAR is driven from WSL with Windows path conversion; the simulation steps cannot run without that host. & \texttt{active\_learning\_pi/\allowbreak al/\allowbreak ecadstar.py} \\
The residual GP is implemented with GPyTorch as an ApproximateGP trained with a VariationalELBO. & \texttt{active\_learning\_pi/\allowbreak al/\allowbreak gp\_error\_surrogate.py} \\
Stale reference document: dated 2026-06-22, describes the exp050 Tier-A configuration, lists six anchors and a heatmap weight of 4.0, neither matching the exp059 dataset nor its config. Usable for method description, not for numbers. & \texttt{docs/\allowbreak normalization-and-losses.md} \\
Stale reference document: written for the exp054-to-exp057 lineage, with its own caveat that exact epoch counts, weights and off-anchor MHz lists are config-specific and should be taken from the experiment's config.yaml. & \texttt{docs/\allowbreak training-procedure.md} \\
Configuration files of record: capacity phase and AL fine-tune, the latter injected at runtime by VAE\_CONFIG\_PATH. & \texttt{experiments/\allowbreak exp059\_capacity\_freq/\allowbreak config.yaml, experiments/\allowbreak exp059\_capacity\_freq/\allowbreak config\_al\_finetune.yaml} \\
\hline
\end{longtable}

\subsection*{Active Learning for scaling the model}
\begin{longtable}{p{0.56\textwidth} p{0.36\textwidth}}
\hline
Statement in the text & Source artifact \\
\hline
\endfirsthead
\hline
Statement in the text & Source artifact \\
\hline
\endhead
Entry point with a CONFIG block: COMMAND = 'full' runs the whole cycle, PROPOSE\_ONLY = True runs generation and inference only. (Entry-point mechanics removed from the body.) & \texttt{pipelines/\allowbreak active\_learning/\allowbreak run.py} \\
The AL pipeline module refuses to run directly and instructs the user to edit the CONFIG block in run.py. & \texttt{active\_learning\_pi/\allowbreak al/\allowbreak pipeline.py} \\
The cycle step count is computed at runtime and is 8 only when fine-tuning and post-fine-tune evaluation are enabled, otherwise 7 or 6. & \texttt{pipelines/\allowbreak active\_learning/\allowbreak run.py} \\
Primary configuration (residual-GP acquisition) with k\_min 3, k\_max 10, k\_sweep\_per\_pool true, candidates\_per\_k 1600, worst\_per\_k 80, total\_ecad\_per\_cycle 640, candidate\_seed 42 with seed = candidate\_seed + iteration*10000 + K, MHz grid and per-K strata quotas, explore\_candidates\_per\_k 160 over 120-320 MHz quantized to 5 MHz with band edges [120, 180, 230, 280, 320], min\_uncertainty\_percentile 40, mhz\_strata\_per\_k set with worst\_mhz\_strata\_per\_k null (auto-scaled to worst\_n). & \texttt{active\_learning\_pi/\allowbreak config/\allowbreak exp059\_gp\_error.json} \\
Alternative configurations: Monte-Carlo acquisition and equal-budget random control. & \texttt{active\_learning\_pi/\allowbreak config/\allowbreak exp059.json, active\_learning\_pi/\allowbreak config/\allowbreak exp059\_random.json} \\
Exp057-era documented budgets conflict with the exp059 configuration: candidates\_per\_k 400, worst\_per\_k 10, 80 ECAD simulations per cycle, off-anchor frequencies 90/265/435/510 MHz. & \texttt{docs/\allowbreak active-learning.md sec.3-4} \\
GP settings: feature path 'pred', target 'p99\_ae', regressor 'svgp', kappa 0.5, n\_fit 4000, n\_inducing 256, svgp\_epochs 80, seed 0, novelty\_weight 0.0, score\_mode 'mu'. Alternative residual targets peak\_px, peak\_soft\_px, peak\_sharp\_px, struct, mae, mse, imp\_mse. & \texttt{active\_learning\_pi/\allowbreak config/\allowbreak exp059\_gp\_error.json, active\_learning\_pi/\allowbreak al/\allowbreak gp\_error\_surrogate.py} \\
predict\_candidates dispatches between gp\_error, random and mc, defaulting to mc when unset; the GP feature path encodes occupancy (encode\_occupancy\_latent), decodes, re-encodes the model's own predicted heatmap and uses that posterior mean. & \texttt{active\_learning\_pi/\allowbreak al/\allowbreak inference\_pool.py} \\
Docstring conflict: states that gp\_error candidates are ranked 'by UCB mu\_e + kappa*sigma\_e'. & \texttt{active\_learning\_pi/\allowbreak al/\allowbreak inference\_pool.py-179} \\
score\_mode 'mu' resolves to the posterior mean only, with sigma unused; the novelty term enters only under 'ucb\_novelty'. & \texttt{active\_learning\_pi/\allowbreak al/\allowbreak gp\_error\_surrogate.py-352} \\
Documentation states the mean-only default explicitly. & \texttt{docs/\allowbreak gp-error-surrogate.md sec.5} \\
The fitted GP artifact is cleared at the start of every per-K acquisition, so each cycle refits. & \texttt{active\_learning\_pi/\allowbreak al/\allowbreak per\_k\_acquire.py-69} \\
Random control assigns uniform scores from np.random.default\_rng(seed + iteration), preserving the ECAD budget exactly. & \texttt{active\_learning\_pi/\allowbreak al/\allowbreak inference\_pool.py, active\_learning\_pi/\allowbreak config/\allowbreak exp059\_random.json} \\
MC-mode parameters retained in the exp059 config: mc\_passes 4, mc\_latent\_passes 4, mc\_decoder\_passes 8, mc\_decoder\_dropout true, score\_metric 'auto\_bad', weights decoder\_dropout\_rce 1.0, latent\_heatmap\_mse 1.0, p99\_var 1.0, p99\_range 0.5, peak\_loc\_spread 0.5, peak\_cv 0.25, prior\_tension 0.35. & \texttt{active\_learning\_pi/\allowbreak config/\allowbreak exp059\_gp\_error.json} \\
The exp059 acquisition A/B (validate\_acquisition\_ab.py) is still an open task and the claim ledger records acq\_ab\_gp\_vs\_mc as UNCERTAIN. & \texttt{docs/\allowbreak cursor-handoff.md sec.3D} \\
Scoring settings: inference\_mode 'layout', pi\_ref\_mhz 200.0, shared\_temp 1.5, heatmap\_grid\_size 64, normalization\_stats\_json pointing at the unbounded training set. & \texttt{active\_learning\_pi/\allowbreak config/\allowbreak exp059\_gp\_error.json} \\
Board mask used at scoring time is a 64x64 boolean array. & \texttt{configs/\allowbreak binary\_mask.npy} \\
Ingested labels are merged into the AL overlay dataset mixed into training with al\_overlay\_sample\_weight 128.0; fine-tune resumes from checkpoints/last\_model.pt with num\_epochs = checkpoint\_epoch + extra\_epochs, extra\_epochs 45; the merged runtime config is written to config\_al\_finetune.runtime.yaml and passed via VAE\_CONFIG\_PATH. & \texttt{datasets/\allowbreak data\_multifreq\_al\_overlay\_exp059, active\_learning\_pi/\allowbreak config/\allowbreak exp059\_gp\_error.json} \\
Fine-tune configuration sets layout\_train\_prob 0.9 and occ\_only\_encode\_prob 0.7, with latent\_distill\_weight 2.5, output\_distill\_weight 2.5, eval\_use\_occ\_only\_layout true and early stopping (patience 3, min delta 0.01, kind 'layout\_cross'). & \texttt{experiments/\allowbreak exp059\_capacity\_freq/\allowbreak config\_al\_finetune.yaml-12} \\
write\_runtime\_finetune\_config copies the AL JSON keys over the values in config\_al\_finetune.yaml, so the executed fine-tune values are 0.65 and 0.4. & \texttt{active\_learning\_pi/\allowbreak al/\allowbreak finetune\_run.py-117} \\
Project documentation states the fine-tune mixing probabilities as 0.9 / 0.7, which is not what runs. & \texttt{docs/\allowbreak cursor-handoff.md sec.2, CLAUDE.md} \\
Reporting gates: decision\_min\_rank\_n 30, decision\_min\_ecad\_n 20, infer\_scope 'ingested'. Per-cycle artifacts: scored\_candidates.json, acquisition\_rank\_quality.json, eval\_off\_anchor\_pre\_finetune.json, eval\_off\_anchor\_post\_finetune.json, eval\_cycle\_summary.json, CYCLE\_EVAL\_REPORT.md, DECISION\_REPORT.md, decision\_ledger.json. & \texttt{active\_learning\_pi/\allowbreak config/\allowbreak exp059\_gp\_error.json, active\_learning\_pi/\allowbreak al/\allowbreak decision\_report.py} \\
The loop has run on exp059 for four iterations; the fourth records acquisition mode gp\_error, K values [3,4,5,6,7,8,9,10], candidates\_per\_k 1600, worst\_per\_k 80, total\_candidates 12800, total\_selected 640. & \texttt{active\_learning\_pi/\allowbreak runs/\allowbreak al\_exp059\_gp\_error\_001/\allowbreak iter\_0004/\allowbreak k\_sweep\_summary.json} \\
\hline
\end{longtable}

